\section{Discussion and Key Findings}
The aggregated data provided deep insights into how RAG systems behave under adversarial semantic conditions. By cross-referencing the experimental configurations, we were able to isolate five notable findings that challenge some assumptions commonly held in NLP security and evaluation methodology.

\subsection{Finding 1: The Parametric Shield Hypothesis (Public vs. Private Data)}
Comparing results between Public data (FastAPI documentation) and Private data (internal security documents) revealed a critical variable for RAG safety: the LLM's parametric knowledge. When evaluated on FastAPI data, the 17B model showed a relatively high Inconsistent rate (9.37\% with prompt and 44.71\% without prompt), because the model could cross-reference the mutated context against knowledge it had picked up during pretraining and spot the inconsistencies.

However, when processing Private Data, the LLM completely lacked this ``parametric shield.'' Without any background knowledge to compare against, the Inconsistent rate dropped sharply to just 2.88\%, while the rate of being fooled (Faithful) climbed significantly to 45.19\%. This finding points to a concerning reality: RAG systems deployed in enterprise environments with internal data are considerably more vulnerable to semantic attacks compared to systems operating on publicly available data.

\subsection{Finding 2: The Inverse Scaling Law of Semantic Poisoning}
One of the fundamental assumptions deeply embedded in the AI field is that scaling laws hold across the board: bigger models with greater capabilities are generally thought to be safer, more resistant to adversarial attacks, and better at catching logical errors. However, comparing the unprotected models (No Prompt) in our study reveals a counterintuitive weakness that we call ``The Inverse Scaling Law of Semantic Poisoning.''

Specifically, when the defensive system prompt was removed, the 17B model's Abstain rate plummeted to just 3.61\%. Meanwhile, the smaller 8B model still managed to maintain an Abstain rate of 42.55\% under the same conditions.

This seemingly paradoxical behavior stems from the 17B model's superior instruction-following and reading comprehension abilities. When faced with complex semantically mutated context, the 17B model tends to process and integrate the flawed logic seamlessly. Without a System Prompt telling it to be skeptical, the model defaults to treating the provided context as absolute truth, leading it to confidently produce Faithful responses (51.44\%). On top of that, its large parametric memory means it frequently detects the errors but then takes it upon itself to correct them (Inconsistent: 44.71\%), which violates the groundedness principle that is central to RAG. The 8B model, on the other hand, lacks the cognitive capacity to smoothly synthesize the contradictory logic in the poisoned context; it gets confused and more often chooses to Abstain (42.55\%). From this, we conclude that more capable LLMs are ironically \textit{more} vulnerable to semantic poisoning when unprotected, because their very capabilities enable a kind of self-deception.

\subsection{Finding 3: Reliability and Consistency of LLM-as-a-Judge}
The reliability of LLM-as-a-Judge methods has long been a topic of interest in automated evaluation research. In this study, rather than using different judge versions for different data types, we chose to use a single AI Judge across all five experimental configurations, covering both Public and Private data.

The cross-validation with human evaluations showed that the AI Judge maintained very high accuracy and consistency. Specifically, Cohen's Kappa scores ranged from Substantial to Almost Perfect agreement ($0.738 \le k \le 0.918$). Notably, on the internal dataset (Private Data), the agreement between the AI Judge and the human evaluator peaked at 0.9180. This shows that a 17B judge model is fully capable of making fine-grained semantic distinctions --- such as telling apart whether the RAG system merely repeated misinformation or used implicit knowledge to ``correct to the right answer'' --- at a level comparable to human judgment, in a fair and objective manner. This finding supports the case for using LLM-as-a-Judge in large-scale security evaluation tasks without needing to constantly retune the judge for each data type.

\subsection{Finding 4: The Hallucination Myth in Semantic Attacks}
A significant portion of NLP security research focuses on measuring and reducing ``hallucinations'' in LLMs. However, our data suggests that this focus is somewhat misplaced when it comes to semantic attacks. Across over 2,000 automated evaluations conducted on two different model sizes, the Hallucination rate never exceeded 0.24\%, and in practice stayed effectively at zero.

When faced with corrupted semantic data, the RAG models did not panic and make up random facts. Instead, their failure mode was distinctly binary: the model either played the role of a passive victim, repeating the false claims verbatim (Faithful), or played the role of a confident rebel, detecting the errors and fixing them using its pretrained knowledge (Inconsistent). This finding resonates with the Factual Accuracy Paradox noted in prior work \parencite{cuconasu2024power}, and suggests that future security research should shift its focus from measuring ``hallucinations'' to measuring ``faithfulness to poison.''

\subsection{Finding 5: Prompt Engineering is a Shield, Not a Cure}
Our evaluation clearly demonstrates the protective role of System Prompts. Deploying a defensive prompt noticeably improved model resilience: the 17B model's Abstain rate was boosted from 3.61\% to 60.10\%, and the 8B model's from 42.55\% to 63.70\%. System prompts essentially act as a first layer of ``immune defense'' in the RAG architecture.

However, despite these improvements, not a single configuration in the study reached the $\ge 90\%$ Abstain threshold. The Binomial tests also decisively rejected the null hypothesis ($H_0$) across all experimental scenarios. This result reinforces the core argument of the study: while Prompt Engineering is a necessary first line of defense, it is not by itself sufficient to fully solve the semantic poisoning problem. Lexical instructions alone simply cannot override the deeper forms of semantic manipulation. This limitation highlights the urgent need for structural and architectural defense mechanisms in future RAG systems, such as approaches combining GraphRAG with Controlled Noise Injection.
